\documentclass[11pt]{article}
\usepackage[a4paper,margin=1in]{geometry}
\usepackage{amsmath,amssymb,mathtools,bm}
\usepackage{enumitem,booktabs,array}
\usepackage{tcolorbox}
\usepackage{hyperref}
\usepackage[protrusion=true,expansion=false]{microtype}
\hypersetup{colorlinks=true,linkcolor=blue,urlcolor=blue,citecolor=blue}
\setlist[itemize]{topsep=2pt,itemsep=2pt}
\setlist[enumerate]{topsep=2pt,itemsep=2pt}
\newtcolorbox{keybox}{colback=blue!3!white,colframe=blue!40!black,boxrule=0.4pt,arc=1mm}
\newtcolorbox{alertbox}{colback=red!3!white,colframe=red!40!black,boxrule=0.4pt,arc=1mm}
\newtcolorbox{answerbox}{colback=green!3!white,colframe=green!40!black,boxrule=0.4pt,arc=1mm}
\newtcolorbox{drillbox}{colback=yellow!5!white,colframe=orange!60!black,boxrule=0.4pt,arc=1mm}
\newcommand{\EV}{\mathbb{E}}
\newcommand{\Var}{\mathrm{Var}}
\newcommand{\Cov}{\mathrm{Cov}}
\newcommand{\blank}[1]{\underline{\hspace{#1}}}

\title{E08-04: Hoberg \& Maksimovic (2022, RFS) \\
\large ``Product Life Cycles in Corporate Finance'' --- Active Recall Workbook}
\author{Empirical Corporate Finance II --- Exam Prep}
\date{\today}
\begin{document}\maketitle

\begin{keybox}
\textbf{Paper in one sentence.} HM use text-based 10-K analysis to classify firms into product life-cycle stages (product development, growth, maturity, decline) and show that firms' investment, R\&D, financing, payout, and M\&A activity vary systematically across stages, providing a quantitative life-cycle framework that complements age-based proxies.

\textbf{Placement.} Extends text-based methods (Hoberg-Phillips) to firm dynamics; standard reference for product life-cycle empirical research in corporate finance.
\end{keybox}

\section*{Round 1: Complete Walk-Through}

\subsection*{1.1 Research question}
Traditional life-cycle measures use firm age or dividend history (e.g., DeAngelo et al.\ 2006). These are crude. Can we use firm disclosures to directly identify life-cycle stages?

\subsection*{1.2 Text-based stage measurement}
\begin{itemize}
\item Train topic models (LDA and related) on 10-K Items 1, 7.
\item Identify topic distributions associated with product-development, growth, maturity, decline stages.
\item Assign each firm-year a share in each stage.
\item Validate: product-development firms have higher R\&D, growth firms more capex, mature firms more payouts.
\end{itemize}

\subsection*{1.3 Data}
\begin{itemize}
\item 10-K filings for all Compustat firms 1996--2018.
\item Cross-referenced with compensation, M\&A, capex, R\&D, payout data.
\end{itemize}

\subsection*{1.4 Headline results}
\begin{itemize}
\item Life-cycle stage shares predict investment, R\&D, M\&A, and payout activity better than firm age.
\item Product-development stage: high R\&D, low capex, external financing dependent.
\item Growth stage: high capex, high M\&A activity.
\item Maturity: high payouts, low R\&D, strong cash flows.
\item Decline: divestitures, shareholder distributions, asset sales.
\end{itemize}

\subsection*{1.5 Mechanism}
\begin{itemize}
\item Text captures evolving product portfolio dynamics.
\item Financial policies align with stage-specific investment/cash-flow profiles.
\item Offers a life-cycle-aware ``frame'' for empirical corporate finance.
\end{itemize}

\subsection*{1.6 Implications}
\begin{itemize}
\item Applied finance should control for text-based life-cycle stage in studies of capital structure, payout, M\&A.
\item Policy on innovation incentives should distinguish development- vs.\ growth-stage firms.
\item Text methods keep extending into firm-level strategy measurement.
\end{itemize}

\subsection*{1.7 Caveats}
\begin{alertbox}
\begin{itemize}
\item Topic-model specification affects stage assignment.
\item 10-K language evolves; cross-period comparability may require adjustment.
\item Private-firm life cycle cannot be measured.
\end{itemize}
\end{alertbox}

\section*{Round 2: Level 1 Fill-in}
\begin{enumerate}
\item Stages: product development / growth / \blank{2cm} / decline.
\item Text source: 10-K Items \blank{1cm}, \blank{1cm}.
\item Method: \blank{2cm} modelling.
\item Stage proxies outperform firm-\blank{1cm} proxies.
\item Growth stage: high capex and \blank{2cm} activity.
\end{enumerate}

\section*{Round 3: Level 2 Prompts}
\begin{enumerate}
\item Describe text-based life-cycle-stage construction.
\item Report main stage-level patterns in investment, R\&D, payout.
\item Discuss why text outperforms age-based proxies.
\item Implications for cross-sectional corporate-finance research.
\item Caveats about topic-model specification.
\end{enumerate}

\section*{Round 4: Minimal Scaffolding}
From a blank page: (i) motivation, (ii) method, (iii) findings, (iv) implications, (v) caveats.

\section*{Round 5: Blank Page}
Essay: product life cycles and corporate-finance policies.

\vspace{6pt}
\begin{drillbox}
\textbf{Speed Drill.} (1) Stages. (2) Text source. (3) Method. (4) Advantage over age. (5) Growth-stage feature.
\end{drillbox}

\section*{Quick Reference \& Answer Key}
\begin{answerbox}
\begin{itemize}
\item \textbf{Stages.} Develop / grow / mature / decline.
\item \textbf{Source.} 10-K Items 1 and 7.
\item \textbf{Method.} Topic models (LDA).
\item \textbf{Use.} Controls for life-cycle in corp-fin research.
\end{itemize}
\end{answerbox}

\begin{keybox}
\textbf{30-second exam pitch.} Hoberg \& Maksimovic (2022) extend the text-based tradition into firm life-cycle analysis. Using topic models on 10-K Items 1 and 7 across all Compustat filings 1996--2018, they classify each firm-year into shares of four product life-cycle stages: product development, growth, maturity, and decline. Stage shares predict corporate policies better than firm age: development-stage firms R\&D-intensive and externally financed, growth-stage firms capex- and M\&A-active, mature firms payout-heavy, decline-stage firms divest and return capital. The paper provides a standard quantitative life-cycle measure for applied research and sharpens cross-sectional tests of capital structure, payout, and investment. It also illustrates the continued payoff of text-based firm-level measurement --- from Hoberg-Phillips 2010 similarity to full firm-strategy classification.
\end{keybox}

\end{document}
